% Generated draft: requires mathematical and source review.
\documentclass[11pt]{article}
\usepackage[margin=25mm]{geometry}
\usepackage{fontspec}
\setmainfont{Latin Modern Roman}
\usepackage{amsmath,amssymb,mathtools,unicode-math}
\setmathfont{Latin Modern Math}
\usepackage{enumitem,needspace,xurl,hyperref}
\setlength{\parindent}{0pt}
\setlength{\parskip}{5pt}
\setlength{\emergencystretch}{3em}
\newcommand{\N}{\mathbb{N}}
\newcommand{\Z}{\mathbb{Z}}
\newcommand{\Q}{\mathbb{Q}}
\newcommand{\R}{\mathbb{R}}
\newcommand{\C}{\mathbb{C}}
\newcommand{\F}{\mathbb{F}}
\newcommand{\A}{\mathbb{A}}
\newcommand{\PP}{\mathbb{P}}
\newcommand{\E}{\mathbb{E}}
\newcommand{\eps}{\varepsilon}
\newcommand{\dd}{\,\mathrm{d}}
\begin{document}
\section*{UnsolvedMath ID: 30001871}
\textbf{Source problem code:} catalogue entry
\subsection{Definitions and mathematical statement}
The supplied record asks two questions. It does not specify the ground field, the convention for Newton--Okounkov bodies of non-big divisors, or the precise meaning of a description by polynomial inequalities. The following gives a conditional formulation in the conventional complex setting and identifies the remaining choices explicitly.

Let \(X\) be a smooth integral projective variety over \(\C\), of dimension \(n\), and let \(L\) be a line bundle on \(X\). Define
\[
\operatorname{vol}_X(L)=\limsup_{m\to\infty}\frac{n!}{m^n}\dim_{\C}H^0(X,L^{\otimes m}).
\]
For an integral divisor \(D\), meaning a divisor with integer coefficients, set \(\operatorname{vol}_X(D)=\operatorname{vol}_X(\mathcal O_X(D))\). On a smooth variety such a divisor is Cartier. No positivity assumption on \(L\) or \(D\) is stated in the supplied questions.

Under the usual real-period convention, a real number \(a\) is a period if it can be expressed as an absolutely convergent integral
\[
a=\int_S f(x_1,\ldots,x_r)\,\dd x_1\cdots\dd x_r,
\]
where \(r\geq 1\), \(f\in\Q(x_1,\ldots,x_r)\), and \(S\subseteq\R^r\) is semialgebraic over \(\Q\). Here semialgebraic over \(\Q\) means describable by a finite Boolean combination of polynomial equalities and inequalities with rational coefficients. This interpretation of ``period'' is conventional but is not explicitly defined in the supplied record.

\textbf{Question (1).} For every smooth projective variety \(X\) and every line bundle \(L\) on \(X\), is \(\operatorname{vol}_X(L)\) a period? In the conditional formulation above, the quantification over varieties is over smooth integral projective varieties over \(\C\); the intended ground-field scope of the source remains unspecified.

For the second question, an admissible flag is a chain
\[
Y_\bullet:\quad X=Y_0\supset Y_1\supset\cdots\supset Y_n=\{p\},
\]
where each \(Y_i\) is irreducible of codimension \(i\) in \(X\), and all \(Y_i\) are nonsingular at \(p\). For a nonzero section \(s\in H^0(X,\mathcal O_X(mD))\), its flag valuation \(\nu_{Y_\bullet}(s)\in\Z_{\geq 0}^n\) records successive orders of vanishing along the flag, dividing out the preceding vanishing before restriction at each step. One complete-linear-series convention defines
\[
\Delta_{Y_\bullet}(D)=\overline{\operatorname{conv}}\left(\bigcup_{\substack{m\geq1\\H^0(X,\mathcal O_X(mD))\neq0}}\left\{\frac{\nu_{Y_\bullet}(s)}m:0\neq s\in H^0(X,\mathcal O_X(mD))\right\}\right)\subseteq\R^n.
\]
Under this convention the body is empty if every positive multiple has no nonzero section. For non-big divisors this construction must be distinguished from limiting-body constructions; the supplied record does not select a convention.

\textbf{Question (2).} For every integral divisor \(D\) on a smooth projective variety \(X\), does there exist an admissible flag \(Y_\bullet\), allowed to depend on \(D\), for which \(\Delta_{Y_\bullet}(D)\) is defined by polynomial inequalities with rational coefficients?

One precise reading of the requested property is that \(\Delta_{Y_\bullet}(D)\) is semialgebraic over \(\Q\). Another possible reading asks for a finite list \(P_1,\ldots,P_N\in\Q[x_1,\ldots,x_n]\) such that
\[
\Delta_{Y_\bullet}(D)=\{x\in\R^n:P_j(x)\geq0\text{ for }1\leq j\leq N\}.
\]
The supplied wording does not specify whether Boolean combinations, strict inequalities, or only simultaneous weak inequalities are permitted. It also does not explicitly restate the ambient variety assumptions in Question (2). Thus these displayed conventions explain possible precise interpretations without resolving the source's missing qualifications.
\subsection{Short English statement}
The source asks whether every line-bundle volume on a smooth projective variety is a period, and whether every integral divisor admits a flag whose Newton--Okounkov body can be described by polynomial inequalities over \(\Q\). The ground field, non-big-divisor convention, and allowed form of inequalities require clarification. The supplied open status is dated curation, not an independently verified conclusion.
\subsection{Sources}
\begin{itemize}
\item Supplied local UnsolvedMath registry record, permanent ID \(30001871\), source problem code \texttt{catalogue entry}, in the collection \emph{Oberwolfach Reports Open Problems}. The record supplies proposed year 2011 and no named proposer.
\item Supplied report citation: \emph{Mini-Workshop: New Developments in Newton-Okounkov Bodies} (2012), \emph{Oberwolfach Reports}, pp.\ 2327--2363, DOI \texttt{10.4171/owr/2011/41}. Supplied source URL and literature-source link, labelled ``original OWR report'': \url{https://doi.org/10.4171/owr/2011/41}. The citation's 2012 date is retained alongside the record's proposed year 2011; their relationship is not explained in the supplied material.
\item Supplied literature-source link, labelled ``volume functions of linear series'': \url{https://arxiv.org/abs/1008.3986}. No authors or further bibliographic details are supplied for this reference.
\item Literature assessment dated exactly \texttt{2026-08-21}: the supplied curation labels the problem open and states that the report records both questions as open, and that known transcendental volume examples and Okounkov-body theory do not resolve either general question. This is supplied dated curation, not independently verified literature assessment. The background expressly requests independent verification of time-sensitive status and cited claims.
\item The record labels its statement \texttt{corrected\_verified} and reports that its mathematical content was checked against the official Oberwolfach report, with notation, context, and question form normalized. This verification claim is supplied provenance; no independent source inspection is asserted here.
\end{itemize}
\end{document}
